% 来源：OEBPS/Text/01_52.xhtml（OPF spine 顺序；仅第一册）
\cleardoublepage
\BookEntry[杰克·韦尔奇《赢》]{杰克·韦尔奇《赢》}

\BookCover{images/书单35赢.jpg}{[美] 杰克·韦尔奇、苏茜·韦尔奇著| 余江、玉书译| 中信出版社2013年版}

\begin{BookSummaryBox}
\textbf{“全球第一CEO”韦尔奇21年领导艺术的精华总结}
\par
\end{BookSummaryBox}

\subsubsection*{这本书谈什么？}

《赢：韦尔奇一生的管理智慧》是韦尔奇结合亲身管理实践和大量鲜活案例，撰写的关于管理理念和方式、方法的著作。他从企业中泛滥而空洞的“使命”说起，谈到每家公司都会面临的日常危机管理，再到预算管理、员工招聘等诸多日常工作，韦尔奇逐一提出了自己独特而富有智慧的解决方案，并与广大读者分享了他作为成功管理者的思维方式。
\par

无论你是财务人员，还是生产线上的技术工人，或者是部门管理者、最基层的普通员工，《赢：韦尔奇一生的管理智慧》都将有利于你更好地了解你所在的工作环境，理解陌生部门的工作，甚至理解公司高层的指令和意图，进而帮助你与同事一道改善工作流程，提高工作效率；同时也完善自己的职业规划，更从容地面对未来的挑战。
\par

\subsubsection*{芒格为什么谈起这本书？}

在斯坦福大学法学院的演讲中，有提问者问芒格：“你讨论过可口可乐的失败。你认为苹果犯了哪些错误呢？”芒格照搬了通用电气CEO杰克·韦尔奇的答案：“我没有足够的能力来回答这个问题。”也是为了说明一个道理：当你不了解，也没有相关才能时，不要害怕说出来。
\par

韦尔奇的代表作就是《赢：韦尔奇一生的管理智慧》，他将自己在工作与生活中如何才能“赢”的智慧倾囊相授。但同时芒格提醒人们，有些人总是很自信地回答他们其实并不了解的问题，就像乱跳舞的蜜蜂，会把整个蜂窝搞得乱哄哄，那是一种巨大的错误。没有人期望你什么都懂。自然，这也是韦尔奇的态度。
\par

\subsubsection*{《赢》作者杰克·韦尔奇如是说}

“从1981年到1995年，我们说我们要做到‘世界上最有竞争力的企业’，成为每个市场的第一名或者第二名——把每个表现不佳、达不到要求的子公司卖掉或者关闭。没有人不明白这个目标是什么意思或意味着什么。”
\par

“在成为领先者之前，成功是发展你自己。当你成为领先者，成功是帮助其他人发展。”
\par

“根据目标进行考核。根据目标提供回报。”
\par

“这些年来我不断学到的教训就是，在许多时候，我太谨慎了。我采取的每个措施几乎都应该早点付诸实施。”
\par

“把计划搞清楚。请一些人来完成它，给他们大量的钱。这向来是成功的不二法门。”
\par

“用各种各样的方式发现人才，我总是认为：你遇到的每个人都是另一次面试。”
\par

“怎样才能找到合适的工作：寻找合适的工作需要时间、实验和耐心。你要学会忍受，这是一个令人讨厌而又浪费时间的过程，有起有伏、反反复复。”
\par

“好的工作和糟糕的工作：
\par

1. 人：你需要找与自己志趣相投的人，在职业生涯中应该越早越好。
\par

2. 机遇：任何新工作都应该让自己感觉有所发展，而不是刚刚够用。
\par

3. 未来：到某些公司工作，就像赢取奥运会的奖牌一样，它将使你在今后的职业生涯中总是与成就和荣誉联系起来（公司声誉）；你所从事的每种职业都是一场赌博，它有可能开拓你未来的发展空间，也有可能缩小你的选择范围（所处产业）。
\par

4. 主导权：在职业生涯中，我们对工作的选择往往都是为了实现其他人的理想。所以在处理工作主导型问题时，唯一的方法就是要对自己诚实，要明白自己是在为哪些人工作。
\par

5. 工作内容：如果一份工作不能让你兴奋——只是由于工作内容乏味——那你就没必要再沉浸其中了。”
\par
